\section{The study: design}
\label{sec:design}

The study asks whether per-channel EMD features add anything over band power, on held-out subjects. It keeps the original features as a reference and adds three feature sets.

\emph{Preprocessing and features.} The original features (F0) use the channels averaged and z-scored. The new sets are computed either on the z-scored channel average (``averaged'') or on each channel separately, with a linear trend removed (``per channel''): F1 is Welch band power \cite{welch1967}, the mean power spectral density in each of five bands (1--4, 4--8, 8--13, 13--30, and 30--45 Hz); F2 is the first four IMFs, each summarized by log variance, mean absolute instantaneous frequency, and histogram entropy (12 features per channel); and F3 is Hilbert band power of all IMFs in the same five bands. Per channel, EMD runs on the signal scaled to unit variance, and the IMFs are scaled back. With channels kept separate, F1 and F3 have 15 features and F2 has 36 on the motor-imagery data sets (25 and 60 on EEGMAT with five channels).

\emph{Classifiers and evaluations.} The study uses shrinkage linear discriminant analysis (LDA), from scikit-learn \cite{pedregosa2011} with Ledoit--Wolf shrinkage \cite{ledoit2004}, and $k$-nearest neighbors (KNN, $k = 5$) with Pearson, Euclidean, and, except for F2, symmetrized Itakura--Saito \cite{itakura1968} distances. F1 and F3 enter as log10 power and F2 as it is, each standardized with statistics of the training windows only, except that the Itakura--Saito distance takes linear power; the F0 features skip that standardization, as in the original pipeline. Before the logarithm or the Itakura--Saito distance, a zero is replaced by its column's smallest positive value, or by 1 if the column has none, using all windows of the data set variant, test windows included, and no labels. In band power, this changes only the 30--45 Hz Hilbert band power of EEG MMI: 16 windows with the channels averaged and, with the channels kept separate, 53 values across 37 windows. In F0, it changes the zero-padded slots of every window with fewer than eight EMD rows. There are three evaluations. (1) The original stratified 70/30 split of windows, over \N{c:split_seeds} seeds. (2) Stratified 5-fold cross-validation within each subject, which asks how well a model trained on a subject's own windows does on that subject. (3) Held-out subjects, by leaving one subject out (BCI IV 2a and EEGMAT) or in 10 subject folds (EEG MMI), with features either as they are or standardized within each subject over all of that subject's windows without labels (for the Itakura--Saito distance, divided by each subject's geometric mean instead). That standardization needs the new subject's whole unlabeled session, so it is named wherever a standardized number appears, and in the text, the number without it is given beside it.

\emph{Statistics.} The summary is the mean over subjects of the per-subject balanced accuracy from out-of-fold predictions, with a 95\% bootstrap interval over subjects (2,000 resamples) \cite{efron1993}. Per-subject permutation tests use \N{c:perm_subject_shuffles} shuffles for LDA and Euclidean KNN (200 for the other distances) with Benjamini--Hochberg correction \cite{benjamini1995}. Paired differences between feature sets use Wilcoxon signed-rank tests \cite{wilcoxon1945} across subjects with Holm's correction \cite{holm1979} over the 27 tests of each data set variant, evaluation, and normalization. The test ranks the differences rounded to ten decimal places, so that differences equal up to floating-point error count as equal: zero differences are dropped, tied ones share their average rank, and a comparison whose rounded differences are all zero gets $p = 1$. The bootstrap intervals and mean differences use the unrounded values. With 9 BCI IV 2a subjects, the smallest possible two-sided Wilcoxon $p$ is 0.0039, which Holm over 27 tests turns into at least 0.105, so no paired BCI IV 2a comparison can pass the correction, and percentile intervals over 9 subjects may run too narrow. For BCI IV 2a, the number of subjects who move in the same direction is given next to each feature comparison that the text reports for that data set alone. The study's main run takes about \N{c:total_minutes} minutes on a laptop CPU, not counting the corrected baseline and the extra checks.
